% Version 13: passages changed by the Toro--Vazquez-Cendon description of the flux split, relative to commit a01174c.
% Each block gives the earlier text as comments (% OLD:), followed by the new text that replaces it.

% ===== sections/01_introduction.tex, version-12 lines 54-54, version-13 lines 54-54
% OLD: versions split the flux into a convective part with flow-velocity eigenvalues and a pressure part carrying the acoustic
versions split the flux into a convective part, whose eigenvalues are the flow velocity or zero, and a pressure part carrying the acoustic

% ===== sections/01_introduction.tex, version-12 lines 115-116, version-13 lines 115-116
% OLD: The discretisation follows a Zha--Bilgen splitting \cite{ZhaBilgen1993}. The explicit step transports all quantities
% OLD: with the flow velocity by the path-conservative Rusanov scheme \cite{Pares2006,Castro2006} with material-speed
The discretisation follows the Toro--V\'azquez-Cend\'on splitting \cite{ToroVazquez2012}, with the mass and internal energy
fluxes evaluated at the new velocity. The explicit step transports all quantities with the flow velocity by the path-conservative Rusanov scheme \cite{Pares2006,Castro2006} with material-speed

% ===== sections/01_introduction.tex, version-12 lines 118-119, version-13 lines 118-119
% OLD: phase consistent. The implicit step carries the pressure, its work on each phase and an incremental flux of the
% OLD: conservative transported quantities. Its central operator is supplemented by Laplace dissipation of pressure and
phase consistent. The implicit step carries the pressure, its work on each phase and the increments of the mass and
internal energy fluxes from the old to the new velocity. Its central operator is supplemented by Laplace dissipation of pressure and

% ===== sections/02_governing.tex, version-12 lines 154-160, version-13 lines 154-159
% OLD: To handle the different time scales of the model we split the fluxes in the manner of Zha and Bilgen
% OLD: \cite{ZhaBilgen1993}, which allows a different time discretisation for each sub-system according to its stiffness: every
% OLD: quantity is transported by the flow velocity in one sub-system, and the pressure, its work and the relaxation terms form
% OLD: the other. The terms of \eqref{eq:general} are separated into the advective part, which transports every component of
% OLD: $\mathbf U$ with the flow velocity, including the non-conservative advection $u\,\partial_x\alpha_1$ of the volume fraction, and the
% OLD: pressure--relaxation part, which contains the pressure flux $\mathbf F^p$, the pressure work of each phase $\mathbf B^p\partial_xu$ and the relaxation terms.
% OLD: Explicitly,
To handle the different time scales of the model we split the fluxes in the manner of Toro and V\'azquez-Cend\'on
\cite{ToroVazquez2012}, which allows a different time discretisation for each sub-system according to its stiffness. The terms
of \eqref{eq:general} are separated into the advective part, which transports the masses, the momentum and the kinetic
energy with the flow velocity and contains the non-conservative advection $u\,\partial_x\alpha_1$ of the volume fraction, and the
pressure--relaxation part, which contains the pressure flux, the internal energy fluxes of the phases with their pressure work
$\mathbf B^p\partial_xu$, the enthalpy flux of the mixture and the relaxation terms. Explicitly,

% ===== sections/02_governing.tex, version-12 lines 163-163, version-13 lines 162-162
% OLD: \mathbf F^c(\mathbf U)=\begin{bmatrix}0\\ m_1u\\ m_2u\\ \rho u^2\\ \varepsilon_1u\\ \varepsilon_2u\\ Eu\end{bmatrix},\quad
\mathbf F^c(\mathbf U)=\begin{bmatrix}0\\ m_1u\\ m_2u\\ \rho u^2\\ 0\\ 0\\ \tfrac12\rho u^3\end{bmatrix},\quad

% ===== sections/02_governing.tex, version-12 lines 165-165, version-13 lines 164-164
% OLD: \mathbf F^p(\mathbf U)=\begin{bmatrix}0\\0\\0\\ \bar p\\0\\0\\ \bar pu\end{bmatrix},\quad
\mathbf F^p(\mathbf U)=\begin{bmatrix}0\\0\\0\\ \bar p\\ \varepsilon_1u\\ \varepsilon_2u\\ \big(E-\tfrac12\rho u^2+\bar p\big)u\end{bmatrix},\quad

% ===== sections/02_governing.tex, version-12 lines 198-202, version-13 lines 197-201
% OLD: 0&0&0&\rho_1\kappa_1&u&0&0\\
% OLD: 0&0&0&\rho_2\kappa_2&0&u&0\\
% OLD: 0&0&0&K^c&0&0&u
% OLD: \end{bmatrix},\qquad
% OLD: K^c=\frac{\sum_k\alpha_k\rho_k\kappa_k/\Gamma_k}{\sum_k\alpha_k/\Gamma_k},
\ast&\ast&0&\rho_1(\kappa_1-\Gamma_1e_1)&0&0&0\\
\ast&0&\ast&\rho_2(\kappa_2-\Gamma_2e_2)&0&0&0\\
\ast&\ast&\ast&K^c&0&0&0
\end{bmatrix},\\
K^c=\frac{\sum_k\alpha_k\rho_k\kappa_k/\Gamma_k-\big(E-\tfrac12\rho u^2\big)}{\sum_k\alpha_k/\Gamma_k},

% ===== sections/02_governing.tex, version-12 lines 206-209, version-13 lines 205-210
% OLD: since the volume fraction, the velocity and the specific internal energies are advected, the phases are compressed by
% OLD: the flow at constant specific internal energy, and $E-\tfrac12\rho u^2$, transported conservatively, determines $p$ at the advected
% OLD: volumes and compressed masses, with $\kappa_k$ and $\Gamma_k$ evaluated at $\rho_k$ and the pressure $p$; for stiffened and Noble--Abel stiffened gases $K^c=\sum_k\alpha_k(p+\gamma_k\pi_k)/(\gamma_k-1)\big/\sum_k(\alpha_k-m_kb_k)/(\gamma_k-1)$.
% OLD: It admits the eigenvalues
since the volume fraction and the velocity are advected, the masses are transported at fixed phase internal energies
$\varepsilon_k$, and the kinetic energy is transported with the mass, so that $E-\tfrac12\rho u^2$ is constant and determines $p$ at the
advected volumes and transported masses; $\kappa_k$, $\Gamma_k$ and $e_k$ are evaluated at the phase state $\rho_k$, $p_k$ in the rows of $p_k$
and at $\rho_k$ and the pressure $p$ in $K^c$. The entries $\ast$ are proportional to $u$. For stiffened and Noble--Abel stiffened gases $\kappa_k-\Gamma_ke_k=\Gamma_k\big(b_k(p_k+\gamma_k\pi_k)/(\gamma_k-1)-q_k\big)$ and
$K^c=\sum_km_k\big(b_k(p+\gamma_k\pi_k)/(\gamma_k-1)-q_k\big)\big/\sum_k(\alpha_k-m_kb_k)/(\gamma_k-1)$; both vanish for ideal and stiffened gases
with $q_k=0$. The last three columns of $\mathbf A^c$ vanish, and the eigenvalues are

% ===== sections/02_governing.tex, version-12 lines 211-211, version-13 lines 212-212
% OLD: \lambda^c_{1,\dots,7}=u .
\lambda^c_{1,\dots,4}=u,\qquad \lambda^c_{5,6,7}=0 .

% ===== sections/02_governing.tex, version-12 lines 214-217, version-13 lines 215-217
% OLD: $\mathbf A^c-u\,\mathbf I$ is nilpotent of rank one, so one eigenvector is missing and the sub-system is weakly hyperbolic, as
% OLD: the pressureless gas dynamics it contains. The characteristic speed is the flow velocity alone. The volume fraction
% OLD: is transported only by this sub-system; in the pressure--relaxation sub-system it changes only through the relaxation
% OLD: terms.
One eigenvector of the eigenvalue $u$ is missing, and the sub-system is weakly hyperbolic, as the pressureless gas dynamics it contains.
The characteristic speeds are the flow velocity and zero. The volume fraction is transported only by this sub-system; in
the pressure--relaxation sub-system it changes only through the relaxation terms.

% ===== sections/02_governing.tex, version-12 lines 232-235, version-13 lines 232-236
% OLD: 0&0&0&\Gamma_1p_1&0&0&0\\
% OLD: 0&0&0&\Gamma_2p_2&0&0&0\\
% OLD: 0&0&0&\bar p\big/\sum_k(\alpha_k/\Gamma_k)&0&0&0
% OLD: \end{bmatrix},
\ast&\ast&0&\rho_1\Gamma_1h_1&u&0&0\\
\ast&0&\ast&\rho_2\Gamma_2h_2&0&u&0\\
\ast&\ast&\ast&K^p&0&0&u
\end{bmatrix},\\
K^p=\frac{E-\tfrac12\rho u^2+\bar p}{\sum_k\alpha_k/\Gamma_k},

% ===== sections/02_governing.tex, version-12 lines 239-242, version-13 lines 240-245
% OLD: with $\tilde{\mathbf S}=(\partial\mathbf Q/\partial\mathbf U)\,\mathbf S$. The fourth row is $\partial_x\bar p/\rho$, $\bar p=\alpha_1p_1+\alpha_2p_2$; each phase keeps its
% OLD: density and is heated by its pressure work, $\partial_tp_k=-\Gamma_kp_k\,\partial_xu$, and the internal energy $E-\tfrac12\rho u^2$ receives $-\bar p\,\partial_xu$, the flux $\partial_x(\bar pu)$ of the total energy
% OLD: minus the work $u\,\partial_x\bar p$ on the kinetic energy, at fixed volumes and phase densities. Only the velocity and the phase pressures are
% OLD: coupled, so the characteristic polynomial of $\mathbf A^p$ is $\lambda^5(\lambda^2-c_{\rm w}^2)$ and the eigenvalues are
with $\tilde{\mathbf S}=(\partial\mathbf Q/\partial\mathbf U)\,\mathbf S$, the coefficients evaluated as in $\mathbf A^c$ and the entries $\ast$ again proportional to $u$. The fourth row is $\partial_x\bar p/\rho$,
$\bar p=\alpha_1p_1+\alpha_2p_2$; each phase keeps its volume and density, and its pressure changes by the transport of $\varepsilon_k$ and by
the pressure work, with the coefficient $\rho_k\Gamma_kh_k$ of $\partial_xu$, $h_k$ the specific enthalpy; the internal energy $E-\tfrac12\rho u^2$ receives
the enthalpy flux of the mixture minus the work $u\,\partial_x\bar p$ on the kinetic energy. The first three rows vanish and the last
column has only its diagonal entry, so the characteristic polynomial of $\mathbf A^p$ is $\lambda^3(\lambda-u)^2(\lambda^2-u\lambda-c_{\rm h}^2)$ and the
eigenvalues are

% ===== sections/02_governing.tex, version-12 lines 244-244, version-13 lines 247-250
% OLD: \lambda^p_{1,\dots,5}=0,\qquad \lambda^p_{6,7}=\pm c_{\rm w},\qquad c_{\rm w}^2=\frac1\rho\sum_k\Gamma_k\alpha_kp_k=\sum_kY_k\,\frac{\Gamma_kp_k}{\rho_k} .
\begin{gathered}
\lambda^p_{1,2,3}=0,\qquad\lambda^p_{4,5}=u,\qquad \lambda^p_{6,7}=\tfrac12\Big(u\pm\sqrt{u^2+4c_{\rm h}^2}\Big),\\
c_{\rm h}^2=\sum_kY_k\Gamma_kh_k .
\end{gathered}

% ===== sections/02_governing.tex, version-12 lines 247-251, version-13 lines 253-261
% OLD: The remaining part of the sound speed, $\sum_kY_k\kappa_k$, comes from the compression of the phases by the mass fluxes of the
% OLD: advective sub-system; the two parts add up to the frozen sound speed, $c_{\rm w}^2+\sum_kY_k\kappa_k=c_f^2$. The
% OLD: speed $c_{\rm w}$ of the pressure work is small in a liquid at moderate pressure, and $c_{\rm w}^2$ is negative for negative phase
% OLD: pressures, as in cavitation; the sub-system is never solved on its own but together with the incremental advective fluxes
% OLD: described below, so that the implicit step carries the whole sound speed, $c_f$ or, with relaxation, that of the active closure.
For $c_{\rm h}^2>0$, $\lambda^p_6>0>\lambda^p_7$ whatever the flow velocity, as for the Euler equations \cite{ToroVazquez2012}. For ideal and
stiffened gases with $q_k=0$, $\Gamma_kh_k=c_k^2$, and the pressure--relaxation sub-system carries the whole frozen sound speed. In
general $c_{\rm h}^2+\sum_kY_k(\kappa_k-\Gamma_ke_k)=c_f^2$, the second part coming from the compression of the phases by the mass fluxes
of the advective sub-system. Through $h_k$, $c_{\rm h}$ depends on the reference energies of the phases, $q_k$ or $\varepsilon_r$: for the water of
Sec.~\ref{sec:res_cav} in its initial state $c_{\rm h}=430$\,m/s against $c_f=1569$\,m/s. Where $c_{\rm h}^2<0$, as in this liquid water below $275$\,K
at $10^5$\,Pa and in the low-pressure initial state of the nitromethane of Sec.~\ref{sec:res_shock}, $\lambda^p_{6,7}$ are complex for $u^2<-4c_{\rm h}^2$
and the sub-system alone is not hyperbolic. It is never solved on its own but together with the incremental mass fluxes
described below, so that the implicit step carries the whole sound speed, $c_f$ or, with relaxation, that of the active
closure.

% ===== sections/02_governing.tex, version-12 lines 256-262, version-13 lines 266-274
% OLD: and the pressure--relaxation sub-system implicitly. So that the implicit step contains the whole acoustic coupling,
% OLD: the advective fluxes are evaluated with the velocity of the new time level, written as the old velocity $u^n$ plus an
% OLD: increment, $Xu^{n+1}=Xu^n+X(u^{n+1}-u^n)$ for $X\in\{m_1,m_2,\varepsilon_1,\varepsilon_2\}$, the momentum and the kinetic
% OLD: energy being carried with the new velocity by the incremental mass flux $\sum_km_k(u^{n+1}-u^n)$: the first part is the
% OLD: explicit advective flux, and the incremental fluxes, which depend on the unknown velocity, are solved together with the
% OLD: pressure--relaxation sub-system in Sec.~\ref{sec:implicit}. Since only the advective eigenvalues enter the explicit stability condition, the time
% OLD: step depends on the flow velocity alone, independently of the sound speed and of the relaxation rates.
and the pressure--relaxation sub-system implicitly. As Chen et al.\ \cite{Chen2026} evaluate the enthalpy flux as $h^n(\rho u)^{n+1}$,
the internal energy fluxes are evaluated with the internal energies at $t^n$ and the velocity at $t^{n+1}$, $\varepsilon^n_ku^{n+1}$, the internal
energy flux of the mixture as their sum and the pressure flux as $\bar p^{\,n+1}u^{n+1}$. So that the implicit step contains the whole
acoustic coupling, the mass fluxes are evaluated in the same way, $m^n_ku^{n+1}$; the momentum and the kinetic energy are carried
with the new velocity by the incremental mass flux $\sum_km_k(u^{n+1}-u^n)$. These fluxes are split into an old-velocity part and an
increment, $Xu^{n+1}=Xu^n+X(u^{n+1}-u^n)$ for $X\in\{m_1,m_2,\varepsilon_1,\varepsilon_2\}$: the first part is evaluated in the explicit step together with the advective fluxes, so that the masses and the internal energies receive the same
numerical dissipation, and the increments, which depend on the unknown velocity, are solved together with the
pressure--relaxation sub-system in Sec.~\ref{sec:implicit}. Since only the flow velocity enters the explicit step, the time step depends
on the flow velocity alone, independently of the sound speed and of the relaxation rates.

% ===== sections/03_method.tex, version-12 lines 8-9, version-13 lines 8-9
% OLD: After the notation, Sec.~\ref{sec:advective} presents the explicit update of the
% OLD: advective sub-system and Sec.~\ref{sec:implicit} the implicit update of the pressure--relaxation sub-system with its
After the notation, Sec.~\ref{sec:advective} presents the explicit
transport and Sec.~\ref{sec:implicit} the implicit update of the pressure--relaxation sub-system with its

% ===== sections/03_method.tex, version-12 lines 18-20, version-13 lines 18-20
% OLD: \caption{Schematic of the semi-implicit scheme. The Zha--Bilgen splitting separates the advective
% OLD: sub-system, advanced explicitly by the path-conservative Rusanov scheme with the material speed, from the
% OLD: pressure--relaxation sub-system, solved implicitly together with the incremental fluxes. Every residual evaluation of the pressure system calls the local closure of each cell, which
\caption{Schematic of the semi-implicit scheme. The Toro--V\'azquez-Cend\'on splitting separates the advective
sub-system, advanced explicitly with the old-velocity part of the internal energy fluxes, from the pressure--relaxation
sub-system, solved implicitly with the increments of the mass and internal energy fluxes. Every residual evaluation of the pressure system calls the local closure of each cell, which

% ===== sections/03_method.tex, version-12 lines 32-32, version-13 lines 32-32
% OLD: \subsection{Advective sub-system}
\subsection{Explicit transport}

% ===== sections/03_method.tex, version-12 lines 34-35, version-13 lines 34-37
% OLD: The advective sub-system \eqref{eq:advective} transports every quantity with the flow velocity. Its $\alpha_1$ row,
% OLD: $u\,\partial_x\alpha_1$, is not the derivative of any flux, and the sub-system is discretised by the path-conservative Rusanov scheme
The explicit step advances the advective sub-system \eqref{eq:advective} together with the old-velocity part of the internal
energy fluxes of Sec.~\ref{sec:split}, $\varepsilon_ku^n$ and $(E-\tfrac12\rho u^2)u^n$, so that every quantity is transported with the flow velocity
by the flux $\mathbf F^u(\mathbf U)=u\,(0,m_1,m_2,\rho u,\varepsilon_1,\varepsilon_2,E)^{\mathsf T}$. The $\alpha_1$ row, $u\,\partial_x\alpha_1$, is not the derivative of
any flux, and the step is discretised by the path-conservative Rusanov scheme

% ===== sections/03_method.tex, version-12 lines 40-40, version-13 lines 42-42
% OLD: \mathbf D^-_{i+1/2}+\mathbf D^+_{i+1/2}=\mathbf F^c(\mathbf U^R)-\mathbf F^c(\mathbf U^L)+\int_0^1\mathbf B^c(\boldsymbol\Psi)\,\partial_\ell\boldsymbol\Psi\,\mathrm d\ell ,
\mathbf D^-_{i+1/2}+\mathbf D^+_{i+1/2}=\mathbf F^u(\mathbf U^R)-\mathbf F^u(\mathbf U^L)+\int_0^1\mathbf B^c(\boldsymbol\Psi)\,\partial_\ell\boldsymbol\Psi\,\mathrm d\ell ,

% ===== sections/03_method.tex, version-12 lines 50-50, version-13 lines 52-52
% OLD: the eigenvalue $u$ of the advective sub-system on the two face states, or its maximum over all faces of the same
the larger flow speed of the two face states, or its maximum over all faces of the same

% ===== sections/03_method.tex, version-12 lines 68-69, version-13 lines 70-81
% OLD: and $G^e$ likewise with $v$. For the row of $\alpha_1$ the path integral is taken along the linear path with the face velocity frozen at the mean
% OLD: $\bar u_{i+1/2}=\tfrac12(u^L+u^R)$, $\int_0^1u\,\partial_\ell\alpha_1\,\mathrm d\ell=\bar u_{i+1/2}\,(\alpha_{1,i+1}-\alpha_{1,i})$, and the fluctuation is split
and $G^e$ likewise with $v$. With the mean face velocity $\bar u_{i+1/2}=\tfrac12(u^L+u^R)$ the speed \eqref{eq:rusanov_speed} is
$a_{i+1/2}=|\bar u_{i+1/2}|+\tfrac12|u^R-u^L|$, and \eqref{eq:rusanov} is identically
\begin{equation}
F^e_{i+1/2}[X]=X^{\uparrow}\,\bar u_{i+1/2}-\tfrac12\big(u^L-u^R\big)^+\big(X^R-X^L\big),
\label{eq:rusanov_upwind}
\end{equation}
with $X^\uparrow=X^L$ for $\bar u_{i+1/2}\ge0$, $X^\uparrow=X^R$ otherwise, and $(\cdot)^+=\max(\cdot,0)$. The first term transports $X$ with the mean face velocity
from its upwind side, the second is a dissipation acting only at compressive faces, $u^L>u^R$. With the maximum $a_\Omega$ of
\eqref{eq:rusanov_speed} over the faces of the same direction, the dissipation coefficient is
$\tfrac12(a_\Omega-|\bar u_{i+1/2}|)-\tfrac14(u^R-u^L)\ge\tfrac12(u^L-u^R)^+$.
For the row of $\alpha_1$ the path integral is taken along the linear path with the face velocity frozen at the mean
$\bar u_{i+1/2}$, $\int_0^1u\,\partial_\ell\alpha_1\,\mathrm d\ell=\bar u_{i+1/2}\,(\alpha_{1,i+1}-\alpha_{1,i})$, and the fluctuation is split

% ===== sections/03_method.tex, version-12 lines 116-116, version-13 lines 128-128
% OLD: The time step is given by the CFL condition of the advective sub-system,
The time step is given by the CFL condition of the explicit transport,

% ===== sections/03_method.tex, version-12 lines 129-129, version-13 lines 141-142
% OLD: All contributions of the pressure--relaxation sub-system are treated implicitly and solved together, within a single step, as follows:
The pressure flux, the pressure work and the relaxation terms, together with the increments of the mass and internal energy
fluxes, are treated implicitly and solved together, within a single step, as follows:

% ===== sections/03_method.tex, version-12 lines 142-142, version-13 lines 155-156
% OLD: With the incremental fluxes of Sec.~\ref{sec:split}, the semi-discrete form of \eqref{eq:pressure_sub} is
With the incremental fluxes of Sec.~\ref{sec:split}, the semi-discrete implicit step, \eqref{eq:pressure_sub} without the old-velocity part of
its internal energy fluxes and with the incremental mass fluxes, is

% ===== sections/03_method.tex, version-12 lines 236-236, version-13 lines 250-250
% OLD: The face sweeps $u_f\tau$ during the step, of which the explicit flux has already swept $\bar u^e_f\tau$, $\bar u^e_f=\bar u_{i+1/2}$ of \eqref{eq:pc_rusanov}; the
The face sweeps $u_f\tau$ during the step, of which the explicit flux has already swept $\bar u^e_f\tau$, $\bar u^e_f=\bar u_{i+1/2}$ of \eqref{eq:rusanov_upwind}; the

% ===== sections/03_method.tex, version-12 lines 243-243, version-13 lines 257-260
% OLD: with $(\cdot)^+=\max(\cdot,0)$ and $(\cdot)^-=\min(\cdot,0)$.
with $(\cdot)^-=\min(\cdot,0)$. Added to the explicit flux \eqref{eq:rusanov_upwind}, the increment gives the flux of the step,
$X^L_f(u_f)^++X^R_f(u_f)^--d_f(X^R_f-X^L_f)$ with $d_f$ the dissipation coefficient of the explicit flux given after \eqref{eq:rusanov_upwind}: $m_k$ and $\varepsilon_k$ are transported by the
face velocity $u_f$ from its upwind side with the face values of the explicit state, $\varepsilon_ku_f$ being the internal energy flux of the
pressure--relaxation sub-system, and the explicit step contributes their common dissipation.

% ===== sections/03_method.tex, version-12 lines 597-597, version-13 lines 614-614
% OLD: phases are compressed at constant specific internal energy, as in the advective sub-system \eqref{eq:advective}.
phases are compressed at constant specific internal energy, $m_k$ and $\varepsilon_k$ being transported by the same linear Rusanov operator with a shared limiter.

% ===== sections/03_method.tex, version-12 lines 680-680, version-13 lines 697-697
% OLD: time linearisation introduced by the splitting, explicit velocity in the advective fluxes and implicit velocity
time linearisation introduced by the splitting, explicit velocity in the transport fluxes and implicit velocity

% ===== sections/03_method.tex, version-12 lines 686-686, version-13 lines 703-703
% OLD: where $\mathcal H$ is the spatial discretisation of the split system, $\mathbf U_E$ the state on which the advective sub-system
where $\mathcal H$ is the spatial discretisation of the split system, $\mathbf U_E$ the state on which the explicit transport

% ===== sections/03_method.tex, version-12 lines 721-721, version-13 lines 738-738
% OLD: \item \emph{Explicit advective sub-system.} The face values are reconstructed on $\mathbf U^e$ by \eqref{eq:shared_limiter} with the cap \eqref{eq:phi_cap}, and every
\item \emph{Explicit transport.} The face values are reconstructed on $\mathbf U^e$ by \eqref{eq:shared_limiter} with the cap \eqref{eq:phi_cap}, and every

% ===== sections/05_conclusions.tex, version-12 lines 29-31, version-13 lines 29-33
% OLD: The flux is split following Zha and Bilgen: all quantities, including the volume fraction, are transported explicitly
% OLD: with the velocity of the explicit state by the path-conservative Rusanov scheme with material-speed dissipation, while
% OLD: the pressure, its work on each phase and an incremental flux are treated implicitly. The central implicit operator is
The flux is split following Toro and V\'azquez-Cend\'on, with the mass and internal energy fluxes evaluated at the new
velocity: all quantities, including the volume fraction, are transported
explicitly with the velocity of the explicit state by the path-conservative Rusanov scheme with material-speed
dissipation, while the pressure, its work on each phase and the increments of the mass and internal energy fluxes from the
old to the new velocity are treated implicitly. The central implicit operator is

% ===== figures/scheme.tex, version-12 lines 25-25, version-13 lines 25-25
% OLD: \node[hdr, below=6mm of fullbox] (split) {Zha--Bilgen splitting};
\node[hdr, below=6mm of fullbox] (split) {Toro--V\'azquez-Cend\'on splitting};

% ===== figures/scheme.tex, version-12 lines 31-31, version-13 lines 31-31
% OLD:    {\fontsize{7pt}{8.4pt}\selectfont path-conservative Rusanov,\\ material speed;\\ $\alpha_1$ advected entirely}};
   {\fontsize{7pt}{8.4pt}\selectfont path-conservative Rusanov,\\ material speed, flux $\mathbf F^u$;\\ $\alpha_1$ advected entirely}};

% ===== figures/scheme.tex, version-12 lines 35-35, version-13 lines 35-35
% OLD:    {\fontsize{7pt}{8.4pt}\selectfont $+$ incremental fluxes $X(u^{n+1}-u^n)$;\\ pressure work $+$ relaxations\\ in one local closure}};
   {\fontsize{7pt}{8.4pt}\selectfont $\varepsilon_ku$, $m_ku$ as increments $X(u^{n+1}-u^n)$;\\ pressure work $+$ relaxations\\ in one local closure}};

% ===== figures/scheme.tex, version-12 lines 49-49, version-13 lines 49-49
% OLD:   {\textbf{1. Explicit advective sub-system} on $\mathbf U^e$, from $\mathbf U^{\rm b}$\\
  {\textbf{1. Explicit transport} on $\mathbf U^e$, from $\mathbf U^{\rm b}$\\

